\documentclass[11pt,a4paper]{article}

\usepackage[T1]{fontenc}
\usepackage[utf8]{inputenc}
\usepackage[english]{babel}
\usepackage{lmodern}
\usepackage{microtype}
\usepackage{amsmath,amssymb,amsthm,mathtools}
\usepackage{geometry}
\usepackage{graphicx}
\usepackage{hyperref}
\usepackage[nameinlink,noabbrev]{cleveref}
\usepackage{enumitem}
\usepackage{xcolor}

\geometry{margin=2.7cm}
\DeclareGraphicsExtensions{.pdf,.png,.jpg}
\graphicspath{{./}{fig/}{figs/}{images/}{img/}}
\newcommand{\includegraphicssafe}[2][]{%
  \IfFileExists{#2}{\includegraphics[#1]{#2}}{\fbox{\texttt{Missing figure: #2}}}%
}

\hypersetup{
  colorlinks=true,
  linkcolor=blue!50!black,
  citecolor=blue!50!black,
  urlcolor=blue!50!black,
  pdftitle={Global Stability in the Two Higgs Doublet Model as a Uniform Coercivity Problem},
  pdfauthor={Thomas Hoffbauer}
}

\theoremstyle{plain}
\newtheorem{theorem}{Theorem}[section]
\newtheorem{proposition}[theorem]{Proposition}
\newtheorem{lemma}[theorem]{Lemma}
\newtheorem{corollary}[theorem]{Corollary}

\theoremstyle{definition}
\newtheorem{definition}[theorem]{Definition}
\newtheorem{example}[theorem]{Example}

\theoremstyle{remark}
\newtheorem{remark}[theorem]{Remark}

\newcommand{\R}{\mathbb{R}}
\newcommand{\B}{\mathbb{B}}
\newcommand{\norm}[1]{\lVert #1 \rVert}
\newcommand{\infk}{\inf_{k\in B^3}}

\title{\textbf{Global Stability in the Two Higgs Doublet Model as a Uniform Coercivity Problem}\\
\large A corrected analytic formulation of the reduced potential}
\author{Thomas Hoffbauer\\Bamberg, Germany}
\date{April 2026}

\begin{document}
\maketitle

\begin{abstract}
A recent formal verification of the reduced stability problem in the two Higgs doublet model (2HDM) showed that a widely cited directional criterion is not sufficient for boundedness from below of the full scalar potential. The underlying issue is structural: once the potential is reduced to the form
\[
V(K_0,k)=K_0J_2(k)+K_0^2J_4(k),
\]
global stability is no longer a purely directional question. It is a uniform coercivity problem coupling the scale variable $K_0>0$ to the direction variable $k\in B^3$. In this note we derive the corrected criterion directly and self-containedly by minimizing in the scale variable. We give a complete proof of the equivalence between global boundedness from below and a uniform domination estimate of the form $J_2(k)^2\le 4c\,J_4(k)$ on the region where $J_2(k)<0$. We also exhibit an explicit family showing why directional nonnegativity of $J_4$ alone is insufficient. The result is best understood as a coercivity statement: the linear drift term must be uniformly controlled by the quadratic curvature term.
\end{abstract}

\section{Introduction}

The question whether a scalar potential is bounded from below is fundamental in classical and quantum field theory. In the two Higgs doublet model (2HDM), this question admits a standard reduction to a scalar scale parameter $K_0>0$ and a direction variable $k\in \R^3$ constrained by $\norm{k}^2\le 1$. In reduced form, the potential becomes
\[
V(K_0,k)=K_0J_2(k)+K_0^2J_4(k),
\]
where $J_2$ is affine and $J_4$ is quadratic in $k$ \cite{Maniatis2006,ToobySmith2026}.

A recent formal analysis has shown that a criterion stated only in terms of directional information is not sufficient for global boundedness from below \cite{ToobySmith2026}. The point is elementary but important: for fixed direction $k$, the behavior of $V(K_0,k)$ depends not only on the sign of $J_4(k)$, but on the quantitative relation between $J_2(k)$ and $J_4(k)$.

The aim of this note is to formulate this correctly in a self-contained analytic way. We show that global stability is equivalent to a uniform coercivity estimate on direction space. In particular, it is not enough that the quartic part be directionally nonnegative; rather, the negative part of the linear term must be uniformly dominated by the quadratic part.

For general background on 2HDMs and orbit-space formulations, see for example \cite{Branco2012,Ivanov2007,Ivanov2008}. The present note is deliberately narrower in scope: it focuses only on the logic of the reduced bounded-from-below problem.

\section{Reduced potential and global stability}

Let
\[
B^3:=\{k\in \R^3:\norm{k}^2\le 1\}.
\]
The reduced 2HDM potential has the form
\begin{equation}\label{eq:reduced-potential}
V(K_0,k)=K_0J_2(k)+K_0^2J_4(k),
\qquad K_0>0,\ k\in B^3,
\end{equation}
where $J_2:B^3\to\R$ is affine and $J_4:B^3\to\R$ is quadratic \cite{Maniatis2006,ToobySmith2026}.

\begin{definition}
The reduced potential $V$ is called \emph{globally stable} if there exists $c\in\R$ such that
\[
V(K_0,k)\ge c
\qquad\text{for all }K_0>0,\ k\in B^3.
\]
\end{definition}

This is a genuinely global condition with quantifier pattern
\[
\exists c\in\R\ \forall k\in B^3\ \forall K_0>0.
\]
Already at this level one sees that a statement involving only directional signs cannot automatically imply global boundedness unless the dependence on $K_0$ has been analyzed uniformly.

For fixed $k$, define
\[
f_k(t):=tJ_2(k)+t^2J_4(k),\qquad t>0.
\]
The stability problem is therefore equivalent to the existence of a uniform lower bound for the family $(f_k)_{k\in B^3}$.

\section{Minimization in the scale variable}

We begin with the elementary one-dimensional minimization problem.

\begin{lemma}\label{lem:radial-infimum}
Fix $k\in B^3$, and write
\[
D:=J_2(k),\qquad Q:=J_4(k).
\]
Then
\[
\inf_{t>0}(tD+t^2Q)=
\begin{cases}
0, & D\ge 0,\ Q\ge 0,\\[1ex]
-\dfrac{D^2}{4Q}, & D<0,\ Q>0,\\[2ex]
-\infty, & Q<0,\\[1ex]
-\infty, & Q=0,\ D<0,\\[1ex]
0, & Q=0,\ D\ge 0.
\end{cases}
\]
\end{lemma}

\begin{proof}
We distinguish cases.

\begin{enumerate}[label=\arabic*.,leftmargin=2em]
\item \textbf{Case $Q<0$.} Then $tD+t^2Q\to -\infty$ as $t\to\infty$, so the infimum is $-\infty$.

\item \textbf{Case $Q=0$.} Then $f_k(t)=tD$. If $D<0$, then $tD\to -\infty$ as $t\to\infty$. If $D\ge 0$, then $\inf_{t>0}tD=0$, approached as $t\downarrow 0$.

\item \textbf{Case $Q>0$.} Completing the square gives
\[
tD+t^2Q
=
Q\left(t+\frac{D}{2Q}\right)^2-\frac{D^2}{4Q}.
\]
If $D<0$, the minimizer $t_*=-D/(2Q)$ lies in $(0,\infty)$, and the infimum equals $-D^2/(4Q)$. If $D\ge 0$, then $t_*\le 0$, so on the domain $t>0$ the infimum is attained only in the limit $t\downarrow 0$, hence equals $0$.
\end{enumerate}
\end{proof}

The preceding lemma motivates the following definition.

\begin{definition}
Define the \emph{effective radial functional}
\[
E(k):=\inf_{K_0>0}V(K_0,k)=\inf_{t>0}(tJ_2(k)+t^2J_4(k)).
\]
\end{definition}

By \cref{lem:radial-infimum} we obtain the explicit formula
\[
E(k)=
\begin{cases}
0, & J_2(k)\ge 0,\ J_4(k)\ge 0,\\[1ex]
-\dfrac{J_2(k)^2}{4J_4(k)}, & J_2(k)<0,\ J_4(k)>0,\\[2ex]
-\infty, & J_4(k)<0,\\[1ex]
-\infty, & J_4(k)=0,\ J_2(k)<0,\\[1ex]
0, & J_4(k)=0,\ J_2(k)\ge 0.
\end{cases}
\]
Hence the global stability problem is exactly the question whether
\[
\infk E(k)>-\infty.
\]

This already isolates the essential obstruction: it is not enough to know that $J_4(k)\ge 0$; one must control the ratio $J_2(k)^2/J_4(k)$ on the set where $J_2(k)<0$.

\section{Corrected global stability criterion}

We now state and prove the main result in self-contained form.

\begin{theorem}[Corrected global stability criterion]\label{thm:main}
The reduced potential
\[
V(K_0,k)=K_0J_2(k)+K_0^2J_4(k)
\]
is globally stable if and only if there exists $c\ge 0$ such that for every $k\in B^3$,
\[
J_4(k)\ge 0,
\]
and whenever $J_2(k)<0$, one has
\[
J_2(k)^2\le 4c\,J_4(k).
\]
\end{theorem}

\begin{proof}
Assume first that $V$ is globally stable. Then there exists $m\in\R$ such that
\[
V(K_0,k)\ge m
\qquad\text{for all }K_0>0,\ k\in B^3.
\]

We first show that $J_4(k)\ge 0$ for all $k\in B^3$. If there existed $k_0$ with $J_4(k_0)<0$, then
\[
V(K_0,k_0)=K_0J_2(k_0)+K_0^2J_4(k_0)\to -\infty
\quad\text{as }K_0\to\infty,
\]
contradicting the lower bound.

Next, suppose $k\in B^3$ satisfies $J_2(k)<0$. Since we already know $J_4(k)\ge 0$, there are two possibilities.

If $J_4(k)=0$, then
\[
V(K_0,k)=K_0J_2(k)\to -\infty
\quad\text{as }K_0\to\infty,
\]
again impossible. Hence $J_4(k)>0$. Therefore \cref{lem:radial-infimum} applies and gives
\[
E(k)=\inf_{K_0>0}V(K_0,k)=-\frac{J_2(k)^2}{4J_4(k)}.
\]
Since $V\ge m$, we have $E(k)\ge m$, hence
\[
-\frac{J_2(k)^2}{4J_4(k)}\ge m.
\]
Let $c:=\max\{0,-m\}$. Then $c\ge 0$ and
\[
\frac{J_2(k)^2}{4J_4(k)}\le c,
\]
that is,
\[
J_2(k)^2\le 4c\,J_4(k).
\]
This proves necessity.

Conversely, assume there exists $c\ge 0$ such that for all $k\in B^3$:
\begin{enumerate}[label=(\roman*),leftmargin=2em]
\item $J_4(k)\ge 0$,
\item if $J_2(k)<0$, then $J_2(k)^2\le 4cJ_4(k)$.
\end{enumerate}
We must prove that $V(K_0,k)\ge -c$ for all $K_0>0$ and $k\in B^3$.

Fix $k\in B^3$.

If $J_2(k)\ge 0$, then since $J_4(k)\ge 0$,
\[
V(K_0,k)=K_0J_2(k)+K_0^2J_4(k)\ge 0\ge -c.
\]

If $J_2(k)<0$, then by assumption $J_4(k)\ge 0$. In fact $J_4(k)>0$, because otherwise $J_2(k)^2\le 4cJ_4(k)=0$ would imply $J_2(k)=0$, contradiction. Hence we may complete the square:
\[
V(K_0,k)
=
J_4(k)\left(K_0+\frac{J_2(k)}{2J_4(k)}\right)^2
-\frac{J_2(k)^2}{4J_4(k)}.
\]
The square term is nonnegative, so
\[
V(K_0,k)\ge -\frac{J_2(k)^2}{4J_4(k)}\ge -c.
\]
Thus for all $k$ and all $K_0>0$,
\[
V(K_0,k)\ge -c.
\]
Hence $V$ is globally stable.
\end{proof}

\section{Why directional positivity alone fails}

The previous theorem shows that the decisive object is the ratio
\[
\frac{J_2(k)^2}{J_4(k)}
\]
on directions with $J_2(k)<0$, not merely the sign of $J_4(k)$. We now make this failure explicit.

\begin{proposition}[Explicit obstruction]\label{prop:counterexample}
There exist continuous functions $J_2,J_4$ on the interval $[0,1]$ such that
\begin{itemize}[leftmargin=1.8em]
\item $J_4(x)\ge 0$ for all $x\in[0,1]$,
\item $J_4(x)>0$ for all $x>0$,
\item but the reduced potential
\[
V(t,x)=tJ_2(x)+t^2J_4(x)
\]
is not globally bounded from below on $t>0$, $x\in[0,1]$.
\end{itemize}
\end{proposition}

\begin{proof}
Take
\[
J_2(x)\equiv -1,\qquad J_4(x)=x.
\]
Then $J_4(x)\ge 0$ for all $x\in[0,1]$, and $J_4(x)>0$ for every $x>0$. However, for each $x>0$,
\[
\inf_{t>0}( -t + xt^2 ) = -\frac{1}{4x}.
\]
As $x\downarrow 0$, this tends to $-\infty$. Hence
\[
\inf_{x\in(0,1]} \inf_{t>0}( -t + xt^2 )=-\infty.
\]
Therefore $V$ is not globally bounded from below.
\end{proof}

This example captures the entire issue. Directional nonnegativity of the quadratic coefficient does not prevent instability if the coefficient degenerates too quickly relative to the linear term.

The same phenomenon can be phrased more generally.

\begin{corollary}\label{cor:sequence-obstruction}
Suppose there exists a sequence $k_n\in B^3$ such that
\[
J_2(k_n)<0,\qquad J_4(k_n)>0,\qquad
\frac{J_2(k_n)^2}{J_4(k_n)}\to\infty.
\]
Then $V$ is not globally stable.
\end{corollary}

\begin{proof}
By \cref{lem:radial-infimum},
\[
E(k_n)=-\frac{J_2(k_n)^2}{4J_4(k_n)}\to -\infty.
\]
Hence $\inf_{k\in B^3}E(k)=-\infty$, so no global lower bound exists.
\end{proof}

\section{Reformulation as uniform coercivity}

The correct structural interpretation is that global stability is a coercivity statement for a coupled scale-direction system.

Set
\[
D(k):=J_2(k),\qquad Q(k):=J_4(k).
\]

\begin{definition}
We say that the pair $(D,Q)$ is \emph{uniformly coercive} on $B^3$ if there exists $c\ge 0$ such that
\begin{enumerate}[label=(\roman*),leftmargin=2em]
\item $Q(k)\ge 0$ for all $k\in B^3$,
\item whenever $D(k)<0$, one has
\[
D(k)^2\le 4c\,Q(k).
\]
\end{enumerate}
\end{definition}

Then \cref{thm:main} becomes the following reformulation.

\begin{corollary}\label{cor:coercivity}
The reduced 2HDM potential is globally stable if and only if the pair $(D,Q)$ is uniformly coercive on $B^3$.
\end{corollary}

This is the correct replacement for a merely directional positivity statement. The role of the two terms is transparent:
\begin{itemize}[leftmargin=1.8em]
\item $D$ is a \emph{drift term}: if negative, it pushes the potential downward along the radial variable;
\item $Q$ is a \emph{curvature term}: it provides the upward quadratic restoration;
\item stability occurs exactly when the downward drift cannot dominate the upward curvature uniformly across directions.
\end{itemize}
Thus the issue is not geometric positivity alone, but quantitative domination.

\section{A finite-dimensional quadratic-form viewpoint}

The structure can also be encoded in a direction-dependent quadratic form. For fixed $k$, define
\[
u(K_0):=\binom{1}{K_0},
\qquad
M(k):=
\begin{pmatrix}
0 & \tfrac12 J_2(k)\\[1ex]
\tfrac12 J_2(k) & J_4(k)
\end{pmatrix}.
\]
Then
\[
u(K_0)^TM(k)\nu(K_0)=K_0J_2(k)+K_0^2J_4(k)=V(K_0,k).
\]

This does not add new content, but it highlights the mechanism: boundedness from below is a semiboundedness property of a family of quadratic forms restricted to the ray $\nu(K_0)$, and failure occurs when the off-diagonal contribution $J_2(k)$ is not uniformly dominated by the diagonal restoring term $J_4(k)$.

This finite-dimensional reformulation should be understood as an alternative bookkeeping device, not as a new theorem.

\section{Heuristic outlook}

The preceding discussion suggests a more general principle that appears far beyond the 2HDM setting. Consider a scale-dependent family
\[
H(t)=tD_1+t^2D_2,
\qquad t>0,
\]
where $D_1$ and $D_2$ are suitable symmetric forms or operators on a Hilbert space. One expects semiboundedness of $H(t)$ to require a relative control of the lower-order contribution by the higher-order one. In the scalar setting treated here, this appears exactly as the estimate
\[
|D(k)|^2\lesssim Q(k).
\]

We stress that this is only a heuristic analogy in the present note. No general infinite-dimensional theorem is claimed here. The purpose of the comparison is simply to situate the reduced 2HDM problem within a familiar pattern: global lower boundedness is typically a domination problem, not a sign problem.

\section{Conclusion}

The reduced 2HDM potential
\[
V(K_0,k)=K_0J_2(k)+K_0^2J_4(k)
\]
makes the central point transparent: global stability is not a purely directional property. It is a uniform statement coupling direction and scale.

The correct analytic criterion is obtained by minimizing in the scale variable. One finds:
\begin{enumerate}[label=\arabic*.,leftmargin=2em]
\item $J_4(k)\ge 0$ is necessary, but not sufficient;
\item the true obstruction is the possible blow-up of $J_2(k)^2/J_4(k)$;
\item global stability is equivalent to the uniform coercivity estimate
\[
J_2(k)^2\le 4cJ_4(k)
\quad\text{on the region }J_2(k)<0.
\]
\end{enumerate}
Thus the corrected principle is:
\begin{quote}
\emph{Global boundedness from below is equivalent to uniform domination of drift by curvature.}
\end{quote}
That is the structural content behind the correction.

\paragraph{Acknowledgment of source context.}
The conceptual starting point for this note is the recent formal verification work identifying an error in the classical 2006 criterion \cite{ToobySmith2026}, together with the original reduced 2HDM formulation of Maniatis, von Manteuffel, Nachtmann, and Nagel \cite{Maniatis2006}.

\bibliographystyle{alpha}
\bibliography{2hdm_stability_revised}

\end{document}
